\begin{figure}[ht]
\centering
\begin{tikzpicture}[x=1cm, y=1cm, font=\small]

% Absolute error of two approximations of e^x on [0,3]:
%  - degree-4 Taylor polynomial (uses 5 coefficients)
%  - [2/2] Pade approximant (uses the same 5 coefficients)
% Both match e^x through x^4. Near the origin the rational form is more
% accurate; the two cross near x = 2 where the target e^x itself is
% growing without bound and neither local form tracks it.
% y = log10(absolute error), mapped y_cm = 5 + log10(err), clamped [0,5.5].
% x = t in [0,3] mapped x_cm = t*2.6. Values from data/pade_vs_taylor.csv.

\draw[->, thick] (-0.2, 0) -- (8.4, 0) node[right] {$x$};
\draw[->, thick] (0, -0.2) -- (0, 5.9) node[above] {abs.\ error (log scale)};

\foreach \t in {0,0.5,1,1.5,2,2.5,3} {
  \pgfmathsetmacro\xc{\t*2.6}
  \draw (\xc, -0.08) -- (\xc, 0.08) node[below, yshift=-0.12cm] {\t};
}
\foreach \p in {0,-1,-2,-3,-4,-5} {
  \pgfmathsetmacro\yc{5+\p}
  \draw (-0.08, \yc) -- (0.08, \yc) node[left, xshift=-0.1cm] {$10^{\p}$};
}

% Taylor deg-4 error |e^x - T4|.
\draw[very thick, engred, mark=*, mark size=1pt] plot coordinates {
  (0.650,0.000) (1.300,1.453) (1.950,2.353) (2.600,2.998) (3.900,3.920)
  (5.200,4.590) (6.500,5.122) (7.800,5.500) };
\node[engred, anchor=south east] at (7.6, 5.10) {Taylor $[4]$};

% [2/2] Pade error |e^x - R22|.
\draw[very thick, engblue, mark=square*, mark size=1pt] plot coordinates {
  (0.650,0.000) (1.300,0.861) (1.950,1.858) (2.600,2.602) (3.900,3.725)
  (5.200,4.590) (6.500,5.290) (7.800,5.500) };
\node[engblue, anchor=north west] at (2.2, 1.55) {Pad\'e $[2/2]$};

% crossover marker
\fill[enggray] (5.2, 4.59) circle (1.4pt);
\node[enggray, anchor=south east] at (5.15, 4.62) {\footnotesize crossover};

\end{tikzpicture}
\caption[Pad\'e versus Taylor approximation error for the exponential]{Absolute
error of two approximations of $e^{x}$ built from the same five
coefficients: the degree-four Taylor polynomial and the $[2/2]$ Pad\'e
approximant. Near the working point the rational form is roughly three
times the more accurate, because it spends its coefficients on the shape
of the function rather than on a single polynomial trend. Both forms
lose accuracy as $x$ leaves the neighbourhood of the expansion point,
and here they cross near $x = 2$, where the target $e^{x}$ is itself
growing steeply and no five-coefficient local form tracks it. On
functions that level off or carry a pole the rational form keeps its
lead well past the point shown, which is the setting where a Pad\'e
approximant earns its cost.}
\label{fig:vol02:ch04:pade-vs-taylor}
\end{figure}
